% GENERATED FILE. Edit canonical lessons, then run npm run generate:books.
% canonical-source-hash: fnv1a64:8881b9ea80d02697
% canonical-lessons: RU-C264-tebe-stoit, RU-C264-prichina-v-tom, RU-C264-review

\chapter{Advice, and Why}
\label{ch:ru-advice-and-why}
\hlchaptermodality{264}

\begin{chapteropening}
\textbf{By the end of this chapter:} \emph{I can give advice with \ru{тебе} \ru{стоит}, \ru{ты} \ru{должен} and \ru{лучше}, and explain why with \ru{причина} \ru{в} \ru{том}, \ru{что} and \ru{поэтому}.}

Complete RU-C264-review: a friend's problem answered with advice at three strengths, then an explanation with \ru{потому} \ru{что} and \ru{поэтому}.
\end{chapteropening}

\section[\texorpdfstring{\ru{тебе} \ru{стоит} … … (tebé stóit … …)}{тебе стоит … … (tebé stóit … …)}]{\ru{тебе} \ru{стоит}, \ru{лучше}, \ru{ты} \ru{должен} — advice at three strengths}

\label{lesson:RU-C264-tebe-stoit}



\begin{quote}
Say \textquotedblleft{}to cost\textquotedblright{} and \textquotedblleft{}to rest\textquotedblright{}. (\emph{Stóit. Otdykhát.}) Then \textquotedblleft{}it seems to me\textquotedblright{}. (\emph{Mne kázhetsya.}) \emph{\ru{Мне}} was \textquotedblleft{}to me\textquotedblright{}; the same shape for \textquotedblleft{}you\textquotedblright{} is \textbf{\ru{тебе}}, \emph{tebé}.
\end{quote}

\begin{grammarlens}[title={Grammar Lens: \ru{тебе} \ru{стоит}, it is worth your while}]
\begin{quote}
\textbf{\ru{Тебе} \ru{стоит} \ru{отдохнуть}.} — \emph{tebé stóit otdokhnút} — You should rest.
\end{quote}

\textbf{\ru{Стоит}}, \textquotedblleft{}it costs\textquotedblright{}, here means \textquotedblleft{}it is worth it\textquotedblright{}: rest would pay you back. The infinitive follows, and the person advised takes the \textquotedblleft{}to you\textquotedblright{} form: \textbf{\ru{тебе}}, or polite \textbf{\ru{вам}}. \emph{\ru{Отдохнуть}} is the perfective partner of \emph{\ru{отдыхать}}: one rest, taken.
\end{grammarlens}

\begin{grammarlens}[title={Grammar Lens: \ru{лучше}, it is better}]
\begin{quote}
\textbf{— \ru{Уже} \ru{поздно}, \ru{а} \ru{метро} \ru{далеко}. — \ru{Тебе} \ru{лучше} \ru{взять} \ru{такси}.} — \emph{uzhé pózdno, a metró dalekó. tebé lúchshe vzyat taksí} — It's late already and the metro is far. You'd better take a taxi.
\end{quote}

\textbf{\ru{Лучше}} — \emph{lúchshe} — \textquotedblleft{}better\textquotedblright{}: one choice against another. \textbf{\ru{Взять}} is the perfective partner of \emph{\ru{брать}}, the pair with two different roots.
\end{grammarlens}

\begin{grammarlens}[title={Grammar Lens: \ru{ты} \ru{должен}, you must}]
\begin{quote}
\textbf{\ru{Ты} \ru{должен} \ru{позвонить} \ru{маме}.} — \emph{ty dólzhen pozvonít máme} — You must ring Mum.
\end{quote}

\textbf{\ru{Должен}} is a duty, and it agrees like the past: \textbf{\ru{должен}} to a man, \textbf{\ru{должна}} (\emph{dolzhná}) to a woman, \textbf{\ru{должны}} (\emph{dolzhný}) to more than one or to \emph{\ru{вы}}. It can sound strong between friends; \textbf{\ru{стоит}} and \textbf{\ru{лучше}} are gentler.
\end{grammarlens}

\begin{cousinweb}
\textbf{\ru{должен}} comes from the Russian noun for \textquotedblleft{}a debt, a duty\textquotedblright{}, which EtymDB 2.1 traces to Proto-Slavic \emph{*dŭlgŭ}. A duty is a debt: \emph{\ru{Я} \ru{должен} \ru{тебе} \ru{сто} \ru{рублей}} still means \textquotedblleft{}I owe you a hundred roubles\textquotedblright{}.
\end{cousinweb}

\subsection*{Your turn}
Say these aloud:

\begin{itemize}
\raggedright
  \item \emph{Tebé stóit otdokhnút.}
  \item \emph{Tebé lúchshe vzyat taksí.}
  \item \emph{Ty dólzhen pozvonít máme}, then to a woman: \emph{Ty dolzhná pozvonít máme.}
  \item a reason and advice: \emph{Uzhé pózdno, poétomu tebé lúchshe vzyat taksí.}
\end{itemize}

\begin{tcolorbox}[breakable,colback=teal!4,colframe=teal!35!black,title={Before you move on}]
The gentlest advice: \textquotedblleft{}it's worth your while\textquotedblright{}? (\emph{\textbf{\ru{Тебе} \ru{стоит}}} …) \textquotedblleft{}You'd better\textquotedblright{}? (\emph{\textbf{\ru{Тебе} \ru{лучше}}} …) The strongest, to a woman? (\emph{\textbf{\ru{Ты} \ru{должна}}} …) And what form of the verb follows all three? (\textbf{The infinitive.})
\end{tcolorbox}
\section[\texorpdfstring{\ru{причина} \ru{в} \ru{том} … (prichína v tom …)}{причина в том … (prichína v tom …)}]{\ru{причина} \ru{в} \ru{том}, \ru{что} … — giving the reason at length}

\label{lesson:RU-C264-prichina-v-tom}



\begin{quote}
Say \textquotedblleft{}you must ring Mum\textquotedblright{}. (\emph{Ty dólzhen pozvonít máme.}) Then \textquotedblleft{}a reason\textquotedblright{}. (\emph{Prichína.}) Advice asks for a reason back.
\end{quote}

\begin{grammarlens}[title={Grammar Lens: \ru{причина} \ru{в} \ru{том}, \ru{что} …}]
\begin{quote}
\textbf{— \ru{Почему} \ru{ты} \ru{не} \ru{позвонил}? — \ru{Причина} \ru{в} \ru{том}, \ru{что} \ru{я} \ru{потерял} \ru{телефон}.} — \emph{pochemú ty ne pozvoníl? prichína v tom, shto ya poteryál telefón} — Why didn't you ring? The reason is that I lost my phone.
\end{quote}

Word for word, \textquotedblleft{}the reason is in that, that …\textquotedblright{}: \textbf{\ru{в} \ru{том}} points ahead, and \textbf{\ru{что}} brings in the explanation, with its comma. It is weightier than \emph{\ru{потому} \ru{что}}.
\end{grammarlens}

\begin{grammarlens}[title={Grammar Lens: an explanation in three steps}]
What happened, then \textbf{\ru{поэтому}}, then what followed:

\begin{quote}
\textbf{\ru{Я} \ru{потерял} \ru{телефон}, \ru{поэтому} \ru{не} \ru{мог} \ru{позвонить}.} — \emph{ya poteryál telefón, poétomu ne mog pozvonít} — I lost my phone, so I couldn't ring.
\end{quote}

\textbf{\ru{Мог}} is the past of \emph{\ru{мочь}}, one of the few whose form for a man has no \emph{\ru{л}}: \emph{\ru{мог}}, but \emph{\ru{могла}}, \emph{\ru{могли}}.
\end{grammarlens}

\begin{culture}
\begin{quote}
\textbf{\ru{Учитель} \ru{говорил} \ru{быстро}, \ru{и} \ru{я} \ru{не} \ru{понимал} \ru{правило}. \ru{Потом} \ru{он} \ru{повторил} \ru{медленно}, \ru{и} \ru{я} \ru{понял}.} — \emph{uchítel govoríl býstro, i ya ne ponimál právilo. potóm on povtoríl médlenno, i ya pónyal} — The teacher was talking fast and I didn't understand the rule. Then he said it again slowly, and I understood.
\end{quote}

A process with no result, then two deeds done: \textbf{\ru{повторил}}, from the partner of \emph{\ru{повторять}}, and \textbf{\ru{понял}}, from \emph{\ru{понять}}, stressed on the front: \emph{pónyal}.
\end{culture}

\begin{cousinweb}
\textbf{\ru{причина}} is \emph{\ru{при}-} with \textbf{\ru{чин}}, by the usual account \textquotedblleft{}order, rank\textquotedblright{}. EtymDB 2.1 lists \emph{\ru{причина}} itself without a parent, and takes \emph{\ru{чин}} through Old East Slavic \emph{činŭ} to Proto-Slavic \emph{*činŭ}. A reason, on this reading, is what set things in motion.
\end{cousinweb}

\subsection*{Your turn}
\begin{itemize}
\raggedright
  \item \emph{Say it:} \emph{Prichína v tom, shto ya poteryál telefón.}
  \item \emph{Say it:} \emph{Ya poteryál telefón, poétomu ne mog pozvonít.}
  \item \emph{Answer:} \emph{Pochemú ty ne pozvoníl?} Begin with \emph{prichína v tom, shto}
\end{itemize}

\begin{tcolorbox}[breakable,colback=teal!4,colframe=teal!35!black,title={Before you move on}]
\textquotedblleft{}The reason is that …\textquotedblright{}? (\emph{\textbf{\ru{Причина} \ru{в} \ru{том}, \ru{что} …}}) \textquotedblleft{}The teacher said it again, and I understood\textquotedblright{}: which aspect? (\textbf{Perfective:} \emph{\textbf{\ru{повторил}}}, \emph{\textbf{\ru{понял}}}.)
\end{tcolorbox}
\section[(review)]{Review — advice, and the reason for it}

\label{lesson:RU-C264-review}



\begin{quote}
Three pieces of advice: \textquotedblleft{}you should rest\textquotedblright{}, \textquotedblleft{}you'd better take a taxi\textquotedblright{}, \textquotedblleft{}you must ring Mum\textquotedblright{}. (\emph{Tebé stóit otdokhnút. Tebé lúchshe vzyat taksí. Ty dólzhen pozvonít máme.})
\end{quote}

\begin{grammarlens}[title={What you've built}]
Two lessons, and Russian can now advise and explain. Advice at three strengths, each with an infinitive: \textbf{\ru{тебе} \ru{стоит}}, \textbf{\ru{лучше}}, \textbf{\ru{ты} \ru{должен}}. And \textbf{\ru{причина} \ru{в} \ru{том}, \ru{что} …}, with the explanation built from what happened and \textbf{\ru{поэтому}}.
\end{grammarlens}

\subsection*{The exchange}
\begin{quote}
\textbf{— \ru{Я} \ru{не} \ru{понимаю} \ru{правило}.} \textbf{— \ru{По}-\ru{моему}, \ru{тебе} \ru{стоит} \ru{спросить} \ru{учителя}. \ru{Он} \ru{поможет}.}
\end{quote}

\emph{Ya ne ponimáyu právilo} — \textquotedblleft{}I don't understand the rule.\textquotedblright{} \emph{Po-móyemu, tebé stóit sprosít uchítelya. On pomózhet} — \textquotedblleft{}I think you should ask the teacher. He'll help.\textquotedblright{} \textbf{\ru{Спросить}} is the perfective partner of \emph{\ru{спрашивать}}.

\subsection*{Your turn}
\begin{itemize}
\raggedright
  \item \emph{Say it:} \emph{Po-móyemu, tebé stóit sprosít uchítelya.}
  \item \emph{Say it:} \emph{Mne kázhetsya, shto tebé lúchshe otdokhnút.}
  \item \emph{Answer:} \emph{Pochemú ty ne pozvoníl?} Begin with \emph{prichína v tom, shto}, then \emph{poétomu}
  \item \emph{Say it:} advice to a friend who is tired, at all three strengths
\end{itemize}

\begin{tcolorbox}[breakable,colback=teal!4,colframe=teal!35!black,title={Before you move on}]
\textquotedblleft{}You must\textquotedblright{}, to a woman? (\emph{\textbf{\ru{Ты} \ru{должна}.}}) \textquotedblleft{}The reason is that …\textquotedblright{}? (\emph{\textbf{\ru{Причина} \ru{в} \ru{том}, \ru{что} …}}) And the gentlest advice? (\emph{\textbf{\ru{Тебе} \ru{стоит}}} …)
\end{tcolorbox}
